% ==========================================================================
%  Paragraph: Exploration Efficiency Analysis
% ==========================================================================

\paragraph{Exploration Efficiency}
Table~\ref{tab:explored_by_k} reports the mean number of explored
(fully expanded) nodes for each algorithm across all $k$-values.
Two key observations emerge.

First, Iterative-$k$A* and Aggregative-$k$A* explore virtually
identical numbers of nodes (${\approx}40{,}000$), differing by less
than $0.3\%$ at every $k$-value. Both counts remain constant as $k$
increases, confirming that heuristic-guided search prunes the same
fraction of the graph regardless of the goal count. Coupled with the
runtime results in Table~\ref{tab:elapsed_by_k}, this reveals that
the $3.9{\times}$ elapsed-time gap between the two algorithms at
$k{=}100$ is \emph{not} caused by differences in search
efficiency, but rather by \emph{per-node computational overhead}:
Aggregative-$k$A* maintains a single open list whose management cost
grows with $k$ (priority-queue operations over an increasingly
populated frontier), whereas Iterative-$k$A* processes each goal
with a compact, reusable data structure whose per-node cost remains
fixed.

Second, $k$Dijkstra explores approximately $4{\times}$ more nodes
(${\approx}162{,}000$) than either heuristic algorithm, confirming
the pruning power of heuristic guidance. Despite this larger search
space, $k$Dijkstra's per-node cost is minimal --- no heuristic
evaluation, no multi-goal bookkeeping --- which explains why its
total runtime (${\approx}6.4$\,s) remains moderate and constant.

Together, the elapsed-time and explored-node analyses provide a
complete picture: Iterative-$k$A* matches the search efficiency of
Aggregative-$k$A* while avoiding its growing per-node overhead,
and achieves $4{\times}$ better pruning than $k$Dijkstra.
